\subsection{A primitive Prym can suffice when the deck action is faithful}

\begin{claim}{A.C2.17}{Proved}
Let $K$ be finitely generated over $\F_2$, and let $C/K$ be a
smooth proper geometrically connected curve of genus two, with
Raynaud theta curve $\Theta$ as in A.C2.5. Suppose that over
$\overline K$ there is a line bundle $L$ of exact order
\[
 n=\ell^a m^b,\qquad
 \ell\ne m\text{ odd primes},\quad a,b\ge1,
\]
such that every multiplication-by-two orbit on
$(\Z/n)^\times$ contains an index $i$ with $L^i\in\Theta$.
Then the geometrically prime-to-two point-section map of $C$
is injective. No condition is needed on characters of smaller
conductor. The conclusion holds after every finite extension,
for forms satisfying this geometric hypothesis, and for every
hyperbolic open subcurve of $C$.

In particular, an exact-order-fifteen theta point suffices,
without any assumption on theta three- or five-torsion.
Every exact theta point of order $3^a5^b$, with $a,b\ge1$,
also suffices. Combining this with A.C2.8, for any odd $n>1$
the simple sufficient condition is
\[
 L\in\Theta\text{ of exact order }n,\qquad
             \langle2,-1\rangle=(\Z/n)^\times.             \tag{PP.1}
\]
These assertions do not assume that $C$ is ordinary.
\end{claim}
\begin{proof}
We first prove a geometric lemma. Let $k$ be algebraically
closed of characteristic two, and let $D/k$ have genus
$n+1$, where $n\ge15$ is odd. Suppose that a cyclic group
$G$ of order $n$ acts freely on $D$ and that
\[
                            f(D)<n-1.                     \tag{PP.2}
\]
Then $D$ has no nonconstant rational function of degree at
most four.

We use the Castelnuovo--Severi inequality in the following
exact form. If $F=F_1F_2$ is a compositum of function
fields of curves over a perfect field, $d_i=[F:F_i]$, and
their genera are $g,g_1,g_2$, then
\[
              g\le d_1g_1+d_2g_2+(d_1-1)(d_2-1).
                                                               \tag{PP.3}
\]
The opened primary author version of
\citet[Proposition 2.1, author p.~3]{A_PrimPoonen2007}
states this without a restriction on the characteristic.
All composita and degrees in its applications are
specified below.

The curve $D$ is not hyperelliptic. If it were, its
hyperelliptic involution would be unique; uniqueness also
follows from (PP.3), since distinct index-two rational
subfields would give $g(D)\le1$. Hence $G$ descends to
its genus-zero quotient. The action there is faithful,
as its kernel lies in the order-two hyperelliptic group.
A generator of an odd cyclic subgroup of
$\operatorname{PGL}_2(k)$ has a fixed point. Its fibre
on $D$ has degree two and is $G$-invariant. Freeness
partitions its geometric points into orbits of size $n$,
with constant multiplicity within each orbit, making
its degree divisible by $n$. This is impossible.

More generally, a $G$-invariant rational subfield
$k(z)\subset k(D)$ of degree $d<n$ is impossible. If
the induced cyclic action on $\mathbf P^1$ is trivial,
every fibre is invariant. If it is nontrivial, a fibre
above a fixed point is invariant. In either case
freeness would force $n\mid d$.

Degree one is excluded by the genus, and degree two by
nonhyperellipticity. A degree-three pencil is unique:
two distinct index-three rational subfields generate
$k(D)$ and would give $g(D)\le4$ by (PP.3).
Uniqueness makes this pencil $G$-invariant, again
impossible. There is no inseparability exception:
over $k$, an inseparable rational function is a square.
If its degree is at most four, a square root has degree
at most two, which has just been excluded.

Suppose therefore that $z$ has degree four. It is
separable. Its rational subfield cannot be
$G$-invariant, so for some $\delta\in G$ the fields
$k(z)$ and $k(\delta z)$ differ. Put
\[
                 F=k(z,\delta z),\qquad e=[k(D):F].
\]
Then $e\mid4$. The value $e=4$ would make the two
rational subfields equal. If $e=1$, (PP.3) gives
$g(D)\le9$, contrary to $g(D)=n+1\ge16$.
Thus $e=2$. The smooth curve $E$ with function field
$F$ has two degree-two rational maps generating its
function field, so (PP.3) gives $g(E)\le1$.
Genus zero would make $D$ hyperelliptic. Consequently
$g(E)=1$. The extension $k(D)/F$ is separable,
because $k(D)/k(z)$ is separable, and has degree two.
Write $\tau$ for its involution.

This bielliptic involution is unique. Two distinct
degree-two elliptic subfields generate $k(D)$, and
(PP.3) with $d_1=d_2=2$, $g_1=g_2=1$ would give
$g(D)\le5$. Hence every automorphism of $D$ commutes
with $\tau$, in particular every element of $G$.

The fixed-point set of $\tau$ is nonempty:
Riemann--Hurwitz gives total different $2n$ for
$D\to E$. It is $G$-stable, and $G$ acts freely on it.
Its cardinal $r$ is a positive multiple of $n$.
Every ramification point of a separable quadratic
extension in characteristic two has different exponent
at least two, since its zeroth and first lower
ramification groups both have order two.
It follows that $2r\le2n$, so $r=n$.
The ramified Deuring--Shafarevich formula now gives
\[
 f(D)-1=2(f(E)-1)+n,\qquad
                   f(D)=2f(E)+n-1\ge n-1,                 \tag{PP.4}
\]
contradicting (PP.2). The exact source is
\citet[Corollary 1.8, pp.~35--36]{A_PrimCrew1984};
both the ramified statement and its proof were opened.

There is also an elementary proof of the needed lower
bound, which makes its geometric content explicit.
The action of $G$ on $E$ is faithful because its
kernel is contained in $\langle\tau\rangle$.
A nonidentity element of $G$ cannot fix a point on
$E$: its odd-order permutation of an at-most-two-point
fibre would be trivial, giving a fixed point on $D$.
Every fixed-point-free elliptic automorphism is a
translation. Indeed, after choosing an origin write
it as $T_b\circ u$ with $u$ fixing the origin. If
$u\ne1$, the nonzero endomorphism $1-u$ is surjective
and produces a fixed point. Thus $G$ acts by an
order-$n$ group of translations.

Take one ramification orbit $P_0,\ldots,P_{n-1}$
and let $b_i$ be its images on $E$.
The differences $b_i-b_0$ are killed by $n$.
Pullback of divisors, whose ramification indices are
two, shows that
\[
                       T_i=n[P_i-P_0],\quad 1\le i<n,
\]
are killed by two. They are linearly independent over
$\F_2$. Otherwise a nonempty relation with
$\epsilon_i\in\{0,1\}$ would give
\[
 \operatorname{div}(v)=
 n\left(\sum_{i=1}^{n-1}\epsilon_iP_i
       -\left(\sum_{i=1}^{n-1}\epsilon_i\right)P_0\right).
\]
Its divisor is $\tau$-invariant, so $\tau(v)/v$ is
constant. The constant has square one and therefore
equals one in characteristic two. Hence $v\in k(E)$
and every valuation at $P_i$ is even. A nonzero
$\epsilon_i$ gives valuation $n$, which is odd.
This contradiction proves $f(D)\ge n-1$
independently of the formula. The geometric lemma is
proved.

Return to the cyclic \'etale cover $D\to C$ associated
with $L$, and let $\sigma$ generate its deck group.
Its genus is $n+1$. Each nonzero character space
$H^1(C,L^i)$ has dimension one. A.C2.5 shows that
the stated theta intersections make every primitive
doubling orbit Frobenius-nilpotent. There are
$\varphi(n)$ primitive characters, so
\[
 f(D)\le f(C)+n-1-\varphi(n)
       \le n+1-\varphi(n)<n-1,                            \tag{PP.5}
\]
since $\varphi(n)\ge8$. The geometric lemma applies.

Set
\[
 a_0=\sigma^{n/\ell},\qquad b_0=\sigma^{n/m},\qquad
          A=(1-a_0)(1-b_0)\in\operatorname{End}(J_D).
                                                               \tag{PP.6}
\]
These automorphisms have distinct odd prime orders
$\ell,m$. On a character of index $i$, the scalar of
$A$ is nonzero exactly when neither $\ell$ nor $m$
divides $i$, equivalently when $i$ is primitive
modulo $n$. It follows that $A$ kills all geometric
two-primary torsion, with the following integral
justification.

Pass from $\Z_2$ to a finite unramified extension
containing all $n$-th roots of unity. Since $n$ is odd,
the cyclic character idempotents act integrally on
the \'etale Tate module. Their summands are free,
and each rank equals the dimension after reduction
modulo two. The Artin--Schreier character description
in A.C2.5 and A.C2.8 makes every primitive summand
zero. Duality may invert the character indices,
but it preserves primitivity.
On each remaining summand one factor of
(PP.6) is exactly zero. Thus $A$ vanishes on
$T_2J_D$ and on $J_D[2^\infty](\overline K)$.
This is not an inference from an arbitrary
endomorphism being zero modulo two.
Equivalently, $\operatorname{im}(A)$ is an abelian
subvariety of two-rank zero, although the
smaller-conductor factors may be ordinary.

If distinct $P,Q\in C(K)$ gave equal section classes,
make a finite extension defining $L,D,\sigma$ and a
chosen lift $P'$. By A.3 there is a lift $Q'$ with
$[P'-Q']$ two-primary. Applying (PP.6) gives
\[
 P'+a_0b_0P'+a_0Q'+b_0Q'
 \ \sim\
 a_0P'+b_0P'+Q'+a_0b_0Q'.                                \tag{PP.7}
\]
The elements $1,a_0,b_0,a_0b_0$ are distinct, because
the two prime-order subgroups intersect trivially.
The free deck action and the distinct base points
make all eight displayed points distinct.
Thus (PP.7) gives a rational function of degree four,
with simple zeros and poles; in particular it is
separable. This contradicts the geometric lemma.
All geometric arguments are over $\overline K$,
while the lifted two-primary difference is over a
finite extension by A.3. No specialization of the
theta hypothesis is needed. Restriction and
functoriality to the compactification give the
stated forms and open-subcurve assertions.

It remains to check the arithmetic corollaries.
The divisor $\Theta$ is symmetric, so a theta point
$L$ supplies both $L$ and $L^{-1}$. Under (PP.1)
each primitive doubling orbit contains $1$ or $-1$,
which proves the requisite orbit condition.
For $n=15$ the primitive orbits are
\[
              \{1,2,4,8\},\qquad \{14,13,11,7\}.
\]
Hence $f(D)\le8$, whereas a bielliptic alternative
would force $f(D)\ge14$. No smaller-order
theta exclusion is necessary.

For $a,b\ge1$,
\[
 \operatorname{ord}_{3^a}(2)=2\cdot3^{a-1},\qquad
 \operatorname{ord}_{5^b}(2)=4\cdot5^{b-1}.
\]
For the first identity, $v_3(2^2-1)=1$, and expansion
of $(1+3^ru)^3-1$ with $3\nmid u$ raises its valuation
from $r$ to $r+1$. The second follows in the same
way from $v_5(2^4-1)=1$ and
$(1+5^ru)^5-1$. Reduction modulo three and five
gives the exact initial orders. The displayed
orders have greatest common divisor two, so two
has order $\varphi(3^a5^b)/2$ modulo $3^a5^b$.
Moreover $-1$ is not a power of two there:
modulo three its exponent would be odd, whereas
modulo five it would be $2\bmod4$.
Thus $\langle2,-1\rangle$ is the full unit group.

Finally, if an odd $n$ has $r$ distinct prime
divisors, the Chinese remainder theorem gives
\[
  (\Z/n)^\times/\bigl((\Z/n)^\times\bigr)^2
                    \simeq(\Z/2)^r.
\]
Two generators cannot span this quotient when
$r>2$. Condition (PP.1) therefore forces at most
two prime divisors. The one-prime case is A.C2.8
and the two-prime case is the theorem just proved.

At degree fifteen the surviving order-three and
order-five characters may have trivial common
kernel. The endomorphism (PP.6) still kills both
and produces a degree-four divisor relation,
explaining the improvement over A.C2.7.
This does not supply a suitable theta point on
every curve: torsion density and A.C2.14 leave
the varying-order existence question open.
The exact arithmetic certificate
\path{code/BC_two_prime_prym_checks.py} and its
JSON output record 289 passing checks, including
261 two-prime orders and eight exact cyclotomic
identities. The proof above establishes the
unbounded assertions.
\end{proof}
